% The `split` listing with arrows drawn between its Arrive and Await lines (listing marks;
% same technique as the arXiv paper's async-bug example, without the \bugYFix hack).
% A plain snippet (not a figure body, not in a sporkfigure: marks need unscaled positions),
% \input by the deck (frame `split_marks`) and by main.tex (fig:split_marks).
%   slides: listing version v on step v+1; RAW arrow from step 2, WAR arrow from step 3
%   paper:  last listing version, both arrows
% Marks must be in normal text (not inside a TikZ node), so the listing is plain text and
% the arrows live in a separate `remember picture, overlay` picture (code_to_tex/README.md, Marks).
\begin{minipage}{116mm}
% Absolute size: \footnotesize differs between the paper (8pt) and slides (10pt).
\fontsize{9}{10.5}\selectfont
\sporkListing{split}
\end{minipage}%
{\sporkFigFont%
\begin{tikzpicture}[remember picture, overlay, rounded corners=2mm, font=\bfseries]
  % Vertical arrow segments run at fixed distances from the listing's left edge (mark `left`),
  % right of the longest code line they pass.
  \coordinate (rawX) at ([xshift=88mm]pic cs:split-left);
  \coordinate (warX) at ([xshift=116mm]pic cs:split-left);
  \foreach \m in {raw_arrive, raw_await, war_arrive, war_await} {
    \coordinate (\m) at ([xshift=1mm]pic cs:split-\m);
  }
  \sporkReveal<2->{
    \draw [spork arrow, draw=sporkRedFg] (raw_arrive) -- (raw_arrive -| rawX)
      -- node[right, text=sporkRedFg] {RAW} (raw_await -| rawX) -- (raw_await);
  }
  \sporkReveal<3->{
    \draw [spork arrow, draw=sporkBlueFg] (war_arrive) -- (war_arrive -| warX)
      -- node[right, text=sporkBlueFg, align=left] {WAR\\(slot\\reuse)} (war_await -| warX) -- (war_await);
  }
\end{tikzpicture}}
